\section*{Hoofdstuk \ref{ch:b3:quantum-model-systems} --- Kwantummechanica voor chemici: modelsystemen}

\begin{solution}{exo:b3:quantum-model-systems:1}
$E_1 = h^2/8m_eL^2$. (a) $L = \qty{1.00}{nm}$: $E_1 = \qty{6.02e-20}{J} =
\qty{0.376}{eV}$; $\Delta E_{12} = 3E_1$ en $\lambda = hc/3E_1 =
\qty{1099}{nm}$, in het nabije infrarood. (b) $L = \qty{0.50}{nm}$: $E_1$ is vier
keer groter, \qty{2.41e-19}{J} $= \qty{1.50}{eV}$, en $\lambda =
\qty{275}{nm}$, in het ultraviolet.
\end{solution}

\begin{solution}{exo:b3:quantum-model-systems:2}
$n_x^2 + n_y^2 + n_z^2 = 3$ (1,1,1): $g = 1$; $6$ (2,1,1 en permutaties):
$g = 3$; $9$ (2,2,1): $g = 3$; $11$ (3,1,1): $g = 3$; $12$ (2,2,2): $g = 1$.
\end{solution}

\begin{solution}{exo:b3:quantum-model-systems:3}
$[\hat x^2,\hat p]f = -\iu\hbar x^2f' + \iu\hbar(x^2f)' = 2\iu\hbar xf$, dus
$[\hat x^2,\hat p] = 2\iu\hbar\hat x$. $[\hat p^2,\hat x]f = -\hbar^2[(xf)'' -
xf''] = -2\hbar^2f'$, dus $[\hat p^2,\hat x] = -2\iu\hbar\hat p$.
\end{solution}

\begin{solution}{exo:b3:quantum-model-systems:4}
$E_J/hc = BJ(J+1)$: 0, 21.19, 63.56 en \qty{127.12}{cm^{-1}}, met $g = 1$,
3, 5, 7. $I = h/(8\pi^2cB) = \qty{2.643e-47}{kg.m^2}$ (met $c$ in
\unit{cm/s}).
\end{solution}

\begin{solution}{exo:b3:quantum-model-systems:5}
$|\psi_n|^2$ is symmetrisch ten opzichte van $L/2$, dus $\langle x\rangle = L/2$. Met
$\sin^2u = \frac12(1 - \cos2u)$ en twee partiële integraties is
$\langle x^2\rangle = L^2\big(\frac13 - \frac{1}{2n^2\pi^2}\big)$. Voor $n = 1$ is
$\langle x^2\rangle = 0.283L^2$, kleiner dan de klassieke $L^2/3$: de
grondtoestand hoopt zich op in het midden. Naarmate $n$ groeit, wordt de
klassieke waarde teruggevonden.
\end{solution}

\begin{solution}{exo:b3:quantum-model-systems:6}
$1\,\unit{cm^{-1}} = \qty{0.011963}{kJ/mol}$. ZPE(\ce{H2}) $= 2200.6 \times
0.011963 = \qty{26.33}{kJ/mol}$, ZPE(\ce{D2}) $= \qty{18.63}{kJ/mol}$; verschil
\qty{7.69}{kJ/mol}. Dezelfde \omterm{def:b3:quantum-model-systems:oscillator}{krachtconstante}, dus $\tilde\omega \propto \mu^{-1/2}$
en de verwachte verhouding is $\sqrt{\mu_D/\mu_H} = 1.4137$ (atoommassa's
2.01410 en 1.00783); de gemeten verhouding is 1.4127, binnen 0.07\,\% (het
kleine verschil komt van de elektronen, die de kernen niet volmaakt
volgen).
\end{solution}

\begin{solution}{exo:b3:quantum-model-systems:7}
$\int_0^Lx^2(L - x)^2\,\dd x = L^5/30$, dus $N = \sqrt{30/L^5}$. Dan is
\[
\langle\psi_1|\phi\rangle = \sqrt{\tfrac2L}\sqrt{\tfrac{30}{L^5}}\int_0^Lx(L -
x)\sin\tfrac{\pi x}{L}\,\dd x = \frac{\sqrt{60}}{L^3}\cdot\frac{4L^3}{\pi^3} =
\frac{4\sqrt{60}}{\pi^3} = 0.99928 .
\]
De parabool is voor 99.86\,\% grondtoestand (het kwadraat van de overlap).
\end{solution}

\begin{solution}{exo:b3:quantum-model-systems:8}
$\langle f|\hat pg\rangle = -\iu\hbar\int f^*g'\,\dd x = -\iu\hbar[f^*g] +
\iu\hbar\int f^{*\prime}g\,\dd x = \int(-\iu\hbar f')^*g\,\dd x = \langle\hat
pf|g\rangle$, omdat de term tussen haken aan de uiteinden nul is. Voor $\hat x\hat p$, met $\hat x$ en $\hat p$ elk \omterm{def:b3:quantum-model-systems:hermitian}{hermitisch}:
$\langle f|\hat x\hat pg\rangle = \langle\hat xf|\hat pg\rangle = \langle\hat p\hat
xf|g\rangle$, en $\hat p\hat x = \hat x\hat p - \iu\hbar \ne \hat x\hat p$: niet
\omterm{def:b3:quantum-model-systems:hermitian}{hermitisch} (de gesymmetriseerde vorm $\frac12(\hat x\hat p + \hat p\hat x)$ is
het wel).
\end{solution}

\begin{solution}{exo:b3:quantum-model-systems:9}
$L = 6 \times 139.7 = \qty{838}{pm}$, $N = 6$: $\lambda = 8m_ecL^2/(7h) =
\qty{331}{nm}$. De absorptie ligt in het ultraviolet: hexatrieen is
kleurloos.
\end{solution}

\begin{solution}{exo:b3:quantum-model-systems:10}
$\kappa = \sqrt{2m(V_0 - E)}/\hbar$ met $V_0 - E = \qty{0.20}{eV}$:
$\kappa_H = \qty{9.82e10}{m^{-1}}$, $\kappa_D = \qty{1.388e11}{m^{-1}}$. De
voorfactoren vallen weg in de verhouding: $T_H/T_D = \eu^{2a(\kappa_D - \kappa_H)} =
\eu^{4.06} = 58$, de exacte waarde: de barrière is voor beide dik
($\kappa a = 4.9$ en 6.9).
\end{solution}

\begin{solution}{exo:b3:quantum-model-systems:11}
Niveaus $E_m = m^2\hbar^2/2m_eR^2$: $m = 0$ bevat twee elektronen, $m = \pm1$
er vier. Het hoogste bezette niveau is $|m| = 1$, het laagste lege $|m| = 2$: $\Delta E =
3\hbar^2/2m_eR^2 = \qty{9.38e-19}{J}$, $\lambda = \qty{212}{nm}$, in het
ultraviolet, waar ook de sterke absorptie van benzeen ligt.
\end{solution}

\begin{solution}{exo:b3:quantum-model-systems:12}
$\langle r\rangle = \frac{4\pi}{\pi a^3}\int_0^\infty r^3\eu^{-2r/a}\dd r =
\frac{4}{a^3}\cdot\frac{3!}{(2/a)^4} = \frac32a$, \qty{79.4}{pm} voor $a = a_0$.
$\langle 1/r\rangle = \frac{4}{a^3}\cdot\frac{1!}{(2/a)^2} = \frac1a$, dus
$\langle V\rangle = -e^2/(4\pi\varepsilon_0a)$. Omdat $E_1 =
-e^2/(8\pi\varepsilon_0a)$ (uit $a = 4\pi\varepsilon_0\hbar^2/\mu e^2$),
is $\langle V\rangle = 2E_1$, en de gemiddelde kinetische energie is $-E_1$.
\end{solution}

\begin{solution}{pb:b3:quantum-model-systems:1}
\textbf{1.} Elk van de $j$ atomen levert één p-orbitaal; de $j - 1$
koolstofatomen en één stikstofatoom leveren elk één elektron en het andere
stikstofatoom zijn vrije elektronenpaar,
want de lading van het kation is over de keten verdeeld: $N = j + 1$, dus 6,
8 en 10.
\textbf{2.} Zes elektronen vullen $n = 1$, 2, 3: HOMO $n = 3$, LUMO $n = 4$.
\textbf{3.} $\Delta E = E_{N/2+1} - E_{N/2} = (N+1)h^2/8mL^2$ en $\lambda =
hc/\Delta E = 8mcL^2/h(N+1)$.
\textbf{4.} $L_0 = 4l$, $6l$, $8l$: 559, 838 en \qty{1118}{pm}.
\textbf{5.} 147, 257 en \qty{374}{nm}.
\textbf{6.} Alle veel te kort (met factoren 3.6, 2.3 en 1.9): de doos van
het model is te klein, want $\lambda \propto L^2$; de elektronen bewegen
over een grotere afstand dan de keten tussen de stikstofatomen.
\textbf{7.} Met de oorsprong in het midden is $\psi_n \propto \cos(n\pi u/L)$ voor
oneven $n$ en $\sin(n\pi u/L)$ voor even $n$ ($u = x - L/2$): even en oneven
functies van $u$.
\textbf{8.} Als $n + n'$ even is, hebben $\psi_n$ en $\psi_{n'}$ dezelfde
pariteit; hun product maal $u$ is oneven en integreert tot nul over een
symmetrisch interval.
\textbf{9.} HOMO $N/2$ en LUMO $N/2 + 1$ hebben $n + n' = N + 1$, oneven: altijd
toegestaan.
\textbf{10.} Tweede kleurstof: HOMO 4, LUMO 5. $4 \to 6$: som even, verboden.
$3 \to 5$: som even, verboden.
\textbf{11.} De volgende toegestane overgang vanuit $n = 4$ is $4 \to 7$, met
een $\Delta E$ die $(49 - 16)/(25 - 16) = 3.67$ keer groter is: bij $\lambda/3.67$,
\qty{164}{nm} voor de tweede kleurstof, diep in het ultraviolet.
\textbf{12.} Hij hangt alleen af van de symmetrie van de potentiaal ten
opzichte van het midden van de keten, en die symmetrie hebben de echte
symmetrische kleurstoffen ook.
\textbf{13.} $E = hc/\lambda = \qty{3.294e-19}{J} = \qty{2.06}{eV}$.
\textbf{14.} $hc\tilde\omega_e = \qty{0.371}{eV}$.
\textbf{15.} $2hcB = \qty{4.21e-22}{J} = \qty{2.63}{meV}$.
\textbf{16.} $kT = \qty{25.7}{meV}$.
\textbf{17.} Alleen rotatieniveaus ($2hcB \ll kT$); vibratie-excitaties
($15kT$) en elektronische excitaties ($80kT$) zijn praktisch niet bezet.
\textbf{18.} Elektronisch: zichtbaar en ultraviolet; vibratie: infrarood;
rotatie: microgolven (en ver infrarood).
\textbf{19.} $L = \sqrt{\lambda h(N+1)/8mc}$ met $N = 8$: $L = \qty{1283}{pm}$.
\textbf{20.} $\delta = (1283 - 838)/2 = \qty{222}{pm}$.
\textbf{21.} $L = 559 + 445 = \qty{1003}{pm}$, $\lambda = \qty{474}{nm}$, 9.5\,\%
onder 524 nm.
\textbf{22.} $L = 1118 + 445 = \qty{1562}{pm}$, $\lambda = \qty{732}{nm}$, 2.9\,\%
boven 711 nm.
\textbf{23.} In de aromatische ringen die de stikstofatomen dragen: het
$\pi$-systeem houdt niet op bij de stikstofatomen, en de effectieve doos
omvat een deel van elke ring.
\textbf{24.} $\delta/l = 222/139.7 = 1.6$: \textbf{de doos reikt ongeveer
\qty{222}{pm}, anderhalve binding, voorbij elk stikstofatoom}, en met die ene
lengte reproduceert het model de drie kleuren binnen 10\,\%.
\end{solution}
